\documentclass{chsmc}
\chsmcsetup{
	year = 2015,
	part = 1,
	type = B,
	show-answers = false,
}
\begin{document}

\section{填空题：本大题共 8 小题，每小题 8 分，共 64 分.}

\begin{enumerate}
	\item 已知函数
		$f(x)=\begin{cases}
			a-x, & x \in [0,3], \\
			a\log_2 x, & x \in (3,+\infty),
		\end{cases}$
		其中 $a$ 为常数，如果 $f(2) < f(4)$，则 $a$ 的取值范围是 \blankline
	\item 已知 $y = f(x) + x^3$ 为偶函数，且 $f(10) = 15$，则 $f(-10)$ 的值为 \blankline
	\item 某房间的室温 $T$（单位：摄氏度）与时间 $t$（单位：小时）的函数关系为：
		$T = a\sin t + b\cos t$，$t \in (0,+\infty)$，其中 $a$, $b$ 为正实数，
		如果该房间的最大温差为 $10$ 摄氏度，则 $a+b$ 的最大值是 \blankline
	\item 设正四棱柱 $ABCD$-$A_1B_1C_1D_1$ 的底面 $ABCD$ 是单位正方形，
		如果二面角 $A_1$-$BD$-$C_1$ 的大小为 $\dfrac{\pi}{3}$，则 $AA_1 =$ \blankline
	\item 已知数列 $\{a_n\}$ 为等差数列，首项与公差均为正数，
		且 $a_2$, $a_5$, $a_9$ 依次成等比数列，则使得
		$a_1+a_2+\cdots+a_k > 100a_1$ 的最小正整数 $k$ 的值是 \blankline
	\item 设 $k$ 为实数，在平面直角坐标系 $xOy$ 中有两个点集
		$A = \{(x,y)\mid x^2+y^2 = 2(x+y)\}$ 和
		$B = \{(x,y)\mid kx-y+k+3 \geqslant 0\}$，
		若 $A \cap B$ 是单元集，则 $k$ 的值为 \blankline
	\item 在平面直角坐标系 $xOy$ 中，$P$ 为椭圆 $\dfrac{y^2}{4} + \dfrac{x^2}{3} = 1$ 上的动点，
		点 $A(1,1)$, $B(0,-1)$，则 $|PA|+|PB|$ 的最大值为 \blankline
	\item 正 $2015$ 边形 $A_1A_2\cdots A_{2015}$ 内接于单位圆 $O$，任取它的两个不同顶点
		$A_i$, $A_j$，则 $\bigl|\overrightarrow{OA_i} + \overrightarrow{OA_j}\bigr| \geqslant 1$
		的概率为 \blankline
\end{enumerate}

\section{解答题：本大题共 3 个小题，满分 56 分. 解答应写出文字说明、证明过程或演算步骤.}

\begin{enumerate}[resume]
	\item (本题满分 16 分) 数列 $\{a_n\}$ 满足 $a_1 = 3$，对任意正整数 $m$, $n$，
		均有 $a_{m+n} = a_m + a_n + 2mn$.
		\begin{enumerate}[(1)]
			\item 求 $\{a_n\}$ 的通项公式;
			\item 如果存在实数 $c$ 使得 $\displaystyle\sum_{i=1}^{k} \frac{1}{a_i} < c$
				对所有正整数 $k$ 都成立，求 $c$ 的取值范围.
		\end{enumerate}
	\item (本题满分 20 分) 设 $a_1$, $a_2$, $a_3$, $a_4$ 为四个有理数，使得
		$\{a_i a_j \mid 1 \leqslant i < j \leqslant 4\} = \left\{-24,-2,-\dfrac{3}{2},-\dfrac{1}{8},1,3\right\}$.
		求 $a_1+a_2+a_3+a_4$ 的值.
	\item (本题满分 20 分) 已知椭圆 $\dfrac{x^2}{a^2} + \dfrac{y^2}{b^2} = 1$（$a > b > 0$）
		的右焦点为 $F(c,0)$，存在经过点 $F$ 的一条直线 $l$ 交椭圆于 $A$, $B$ 两点，
		使得 $OA \perp OB$，求该椭圆的离心率 $e = \dfrac{c}{a}$ 的取值范围.
\end{enumerate}

\end{document}
